% Options for packages loaded elsewhere
\PassOptionsToPackage{unicode}{hyperref}
\PassOptionsToPackage{hyphens}{url}
\PassOptionsToPackage{dvipsnames,svgnames,x11names}{xcolor}
%
\documentclass[
  letterpaper,
  DIV=11,
  numbers=noendperiod]{scrartcl}

\usepackage{amsmath,amssymb}
\usepackage{iftex}
\ifPDFTeX
  \usepackage[T1]{fontenc}
  \usepackage[utf8]{inputenc}
  \usepackage{textcomp} % provide euro and other symbols
\else % if luatex or xetex
  \usepackage{unicode-math}
  \defaultfontfeatures{Scale=MatchLowercase}
  \defaultfontfeatures[\rmfamily]{Ligatures=TeX,Scale=1}
\fi
\usepackage{lmodern}
\ifPDFTeX\else  
    % xetex/luatex font selection
\fi
% Use upquote if available, for straight quotes in verbatim environments
\IfFileExists{upquote.sty}{\usepackage{upquote}}{}
\IfFileExists{microtype.sty}{% use microtype if available
  \usepackage[]{microtype}
  \UseMicrotypeSet[protrusion]{basicmath} % disable protrusion for tt fonts
}{}
\makeatletter
\@ifundefined{KOMAClassName}{% if non-KOMA class
  \IfFileExists{parskip.sty}{%
    \usepackage{parskip}
  }{% else
    \setlength{\parindent}{0pt}
    \setlength{\parskip}{6pt plus 2pt minus 1pt}}
}{% if KOMA class
  \KOMAoptions{parskip=half}}
\makeatother
\usepackage{xcolor}
\setlength{\emergencystretch}{3em} % prevent overfull lines
\setcounter{secnumdepth}{-\maxdimen} % remove section numbering
% Make \paragraph and \subparagraph free-standing
\makeatletter
\ifx\paragraph\undefined\else
  \let\oldparagraph\paragraph
  \renewcommand{\paragraph}{
    \@ifstar
      \xxxParagraphStar
      \xxxParagraphNoStar
  }
  \newcommand{\xxxParagraphStar}[1]{\oldparagraph*{#1}\mbox{}}
  \newcommand{\xxxParagraphNoStar}[1]{\oldparagraph{#1}\mbox{}}
\fi
\ifx\subparagraph\undefined\else
  \let\oldsubparagraph\subparagraph
  \renewcommand{\subparagraph}{
    \@ifstar
      \xxxSubParagraphStar
      \xxxSubParagraphNoStar
  }
  \newcommand{\xxxSubParagraphStar}[1]{\oldsubparagraph*{#1}\mbox{}}
  \newcommand{\xxxSubParagraphNoStar}[1]{\oldsubparagraph{#1}\mbox{}}
\fi
\makeatother


\providecommand{\tightlist}{%
  \setlength{\itemsep}{0pt}\setlength{\parskip}{0pt}}\usepackage{longtable,booktabs,array}
\usepackage{calc} % for calculating minipage widths
% Correct order of tables after \paragraph or \subparagraph
\usepackage{etoolbox}
\makeatletter
\patchcmd\longtable{\par}{\if@noskipsec\mbox{}\fi\par}{}{}
\makeatother
% Allow footnotes in longtable head/foot
\IfFileExists{footnotehyper.sty}{\usepackage{footnotehyper}}{\usepackage{footnote}}
\makesavenoteenv{longtable}
\usepackage{graphicx}
\makeatletter
\def\maxwidth{\ifdim\Gin@nat@width>\linewidth\linewidth\else\Gin@nat@width\fi}
\def\maxheight{\ifdim\Gin@nat@height>\textheight\textheight\else\Gin@nat@height\fi}
\makeatother
% Scale images if necessary, so that they will not overflow the page
% margins by default, and it is still possible to overwrite the defaults
% using explicit options in \includegraphics[width, height, ...]{}
\setkeys{Gin}{width=\maxwidth,height=\maxheight,keepaspectratio}
% Set default figure placement to htbp
\makeatletter
\def\fps@figure{htbp}
\makeatother

\KOMAoption{captions}{tableheading}
\makeatletter
\@ifpackageloaded{caption}{}{\usepackage{caption}}
\AtBeginDocument{%
\ifdefined\contentsname
  \renewcommand*\contentsname{Table of contents}
\else
  \newcommand\contentsname{Table of contents}
\fi
\ifdefined\listfigurename
  \renewcommand*\listfigurename{List of Figures}
\else
  \newcommand\listfigurename{List of Figures}
\fi
\ifdefined\listtablename
  \renewcommand*\listtablename{List of Tables}
\else
  \newcommand\listtablename{List of Tables}
\fi
\ifdefined\figurename
  \renewcommand*\figurename{Figure}
\else
  \newcommand\figurename{Figure}
\fi
\ifdefined\tablename
  \renewcommand*\tablename{Table}
\else
  \newcommand\tablename{Table}
\fi
}
\@ifpackageloaded{float}{}{\usepackage{float}}
\floatstyle{ruled}
\@ifundefined{c@chapter}{\newfloat{codelisting}{h}{lop}}{\newfloat{codelisting}{h}{lop}[chapter]}
\floatname{codelisting}{Listing}
\newcommand*\listoflistings{\listof{codelisting}{List of Listings}}
\makeatother
\makeatletter
\makeatother
\makeatletter
\@ifpackageloaded{caption}{}{\usepackage{caption}}
\@ifpackageloaded{subcaption}{}{\usepackage{subcaption}}
\makeatother

\ifLuaTeX
  \usepackage{selnolig}  % disable illegal ligatures
\fi
\usepackage{bookmark}

\IfFileExists{xurl.sty}{\usepackage{xurl}}{} % add URL line breaks if available
\urlstyle{same} % disable monospaced font for URLs
\hypersetup{
  colorlinks=true,
  linkcolor={blue},
  filecolor={Maroon},
  citecolor={Blue},
  urlcolor={Blue},
  pdfcreator={LaTeX via pandoc}}


\author{}
\date{}

\begin{document}


\section{Measuring Consumer Preferences for Used Battery Electric
Vehicles}\label{measuring-consumer-preferences-for-used-battery-electric-vehicles}

\textbf{Zain Hoda}\textsuperscript{1}, \textbf{Xiatian
Iogansen}\textsuperscript{1}, \textbf{John Helveston}\textsuperscript{1}

\textsuperscript{1}Department of Engineering Management and Systems
Engineering, George Washington University, Washington, D.C., USA

\begin{center}\rule{0.5\linewidth}{0.5pt}\end{center}

\subsection{1 Motivation and Research
Questions}\label{motivation-and-research-questions}

Used vehicles account for approximately 70\% of all vehicle transactions
in the United States {[}1{]}, yet the used battery electric vehicle
(BEV) segment has received comparatively little systematic inquiry from
researchers or policymakers. This inattention is increasingly difficult
to justify: used EV sales rose 34.2\% year-over-year as of February
2025, while average days' supply declined 21.5\% over the same period
{[}2{]}, signaling demand growth that is outpacing inventory. The used
BEV market also carries significance beyond its transaction volume. It
extends the productive life of vehicles and battery packs, supports
circular economy objectives {[}3{]}, and reduces the effective entry
price for lower- and middle-income households seeking the operating cost
advantages of electric propulsion without the capital outlay of a new
vehicle purchase.

Consumer demand for used BEVs remains constrained, however, by barriers
that do not apply equally to conventional used vehicles. Survey evidence
indicates that 52\% of U.S. adults are unlikely to consider a used EV
purchase, compared with 28\% who express openness to doing so {[}4{]}.
Used BEVs share certain adoption barriers with new EVs --- range
anxiety, charging infrastructure availability --- but carry additional
purchase risks specific to the secondary market. Among these, battery
condition is the most consequential: it is the primary determinant of a
used BEV's residual range, longevity, and resale value, yet it is
largely unobservable to prospective buyers at the point of sale. Unlike
mileage or model year, which serve as widely accepted proxies for
mechanical wear in conventional vehicles, battery degradation
trajectories vary substantially among BEVs of identical vintage and
odometer reading as a function of prior charging behavior, thermal
exposure, and duty cycle {[}5{]}. This structural information asymmetry
creates conditions under which buyers systematically discount older BEVs
as a hedge against unobserved deterioration, potentially suppressing
market prices below the true value of well-maintained vehicles and
dampening demand across the segment.

This study addresses three interrelated research questions:

\begin{enumerate}
\def\labelenumi{\arabic{enumi}.}
\tightlist
\item
  What is consumers' willingness to pay (WTP) for a used BEV powertrain
  relative to a comparable internal combustion engine (ICE) vehicle, and
  how does this WTP vary across price budget tiers?
\item
  How do the principal vehicle attributes --- driving range, vehicle
  age, and powertrain type --- jointly determine WTP for used BEVs?
\item
  To what extent do preference structures and WTP magnitudes differ
  between consumers who self-identify as likely versus unlikely future
  BEV adopters?
\end{enumerate}

\begin{center}\rule{0.5\linewidth}{0.5pt}\end{center}

\subsection{2 Background and
Contribution}\label{background-and-contribution}

The EV adoption literature is dominated by studies of new vehicle
buyers, in which purchase price, operating costs, driving range,
recharging time, and charging infrastructure availability are
consistently identified as the primary determinants of choice {[}6{]}.
Relatively little of this work extends to the used-vehicle segment,
which serves a distinct population with different cost constraints,
vehicle quality expectations, and information environments.
Affordability, vehicle condition, and reliability are the dominant
evaluation criteria for used-vehicle buyers {[}6{]}; total cost of
ownership and fuel costs are also salient, particularly among
cost-constrained households. Average used-vehicle transaction prices
have remained elevated since the post-pandemic period, reaching \$25,547
in April 2025 {[}7{]}, compressing the potential cost advantage of used
BEVs for budget-sensitive buyers.

Among the limited studies that do examine the used EV market
specifically, qualitative and descriptive analyses predominate.
Segmentation work identifies three broadly distinct prospective buyer
types: those for whom purchase cost constitutes a binding constraint,
those resistant to the behavioral and technological adjustments
associated with electric propulsion, and those motivated by a favorable
cost-to-benefit ratio relative to accelerated new EV depreciation
{[}8{]}. The demographic profile of observed used EV owners ---
disproportionately male, middle-aged, higher-income, and highly educated
--- closely mirrors that of new EV owners {[}9{]}, indicating that the
secondary market has not yet substantially extended access to
lower-income households. Pre-purchase concerns center on battery
condition, charging access, and vehicle price; lower-income prospective
buyers weight both price and charging availability more heavily than
higher-income counterparts {[}10{]}.

This literature, while informative about the qualitative structure of
buyer heterogeneity, has not produced disaggregated quantitative
estimates of WTP for powertrain type in the used BEV market. The present
study addresses that gap through a discrete choice experiment (DCE)
administered to a nationally distributed sample of prospective
used-vehicle buyers in the United States. The DCE design recovers
implicit monetary valuations for vehicle attributes directly from
observed choices between competing alternatives, enabling WTP estimation
disaggregated by price budget tier, adoption propensity, vehicle type
(car vs.~SUV), driving range, and vehicle age. To the authors'
knowledge, this constitutes the first systematic quantification of used
BEV willingness to pay at this level of segmentation in a U.S. context.

\begin{center}\rule{0.5\linewidth}{0.5pt}\end{center}

\subsection{3 Data and Methods}\label{data-and-methods}

\subsubsection{3.1 Survey Design}\label{survey-design}

The survey targets U.S. residents aged 18 years or older who report an
intention to purchase a used car or SUV within the following two years.
The core instrument is a DCE in which respondents complete six repeated
choice tasks. Each task presents three used vehicle alternatives ---
differentiated on observed attributes --- alongside a ``none of the
above'' opt-out option. Respondents select their preferred alternative
in each task. This repeated-measures design exploits within-respondent
variation across attribute combinations to estimate preference
parameters more precisely than a single-task design would permit.

Vehicle attributes included in the DCE are mileage, powertrain type,
operating cost, model year, and purchase price. Powertrain is defined at
three levels: conventional internal combustion engine vehicle (ICEV),
gas hybrid electric vehicle (HEV), and battery electric vehicle (BEV),
with the ICEV serving as the reference category throughout. Operating
cost is expressed in cents per mile driven. Purchase price levels are
segmented by vehicle type and consumer budget tier to maintain
contextual plausibility of the attribute ranges: low-budget cars
(\$10,000--\$20,000), high-budget cars (\$20,000--\$40,000), low-budget
SUVs (\$15,000--\$25,000), and high-budget SUVs (\$25,000--\$45,000).
All attribute definitions are presented to respondents before the choice
tasks begin to ensure consistent interpretation across the sample.

In addition to the DCE, the survey collects information on respondents'
current vehicle ownership and primary vehicle characteristics, future
purchase intentions and anticipated budget, EV-specific knowledge and
attitudes (assessed through 5-point Likert scales covering range
anxiety, charging infrastructure concerns, environmental beliefs, cost
perceptions, and battery-related concerns), and standard
sociodemographic characteristics including household income, housing
type, education level, and employment status.

\subsubsection{3.2 Data Collection and
Sample}\label{data-collection-and-sample}

The survey was administered online between December 2025 and April 2026
across two established research platforms. Respondents who failed
embedded attention checks, exhibited within-survey response
inconsistencies, or provided uninformative filler text in open-ended
fields were removed from analysis. The final analytic sample comprises
\textbf{2,916 respondents} distributed across the continental United
States. Median survey completion time was approximately 14 minutes.

\subsubsection{3.3 Estimation}\label{estimation}

Consumer choice is modeled within the random utility framework. Each
respondent \(i\) is assumed to select the alternative \(j\) that
maximizes their latent utility \(U_{ij}\) in each choice task, where
utility is a function of observed vehicle attributes and an unobserved
idiosyncratic error term assumed to follow an independent Gumbel
distribution.

A multinomial logit (MNL) model, estimated via maximum likelihood,
provides the baseline specification under the assumption of preference
homogeneity across respondents. To relax this restriction --- which is
inconsistent with the substantial heterogeneity in EV attitudes
documented in the literature --- the primary analysis employs a mixed
logit (MXL) model. In the MXL specification, each of the eight
preference parameters is treated as a draw from an independent normal
mixing distribution:

\[\beta_{ik} \sim \mathcal{N}(\mu_k,\, \sigma_k^2)\]

where \(\mu_k\) is the population mean valuation for attribute \(k\) and
\(\sigma_k\) captures the dispersion of that valuation across
respondents. A statistically significant estimate of \(\sigma_k\)
constitutes evidence that preferences for attribute \(k\) are genuinely
heterogeneous in the population, rather than uniform. The eight
parameters specified as randomly distributed are the BEV powertrain
indicator, HEV powertrain indicator, BEV driving range, vehicle mileage,
vehicle age, operating cost, purchase price, and the
alternative-specific constant for the no-choice option.

Six MXL models are estimated, stratified by self-reported adoption
propensity (likely vs.~unlikely BEV adopter) and vehicle segment (full
sample, car segment, SUV segment), yielding models for all combinations
of these two dimensions. All models are estimated in R using the
\texttt{logitr} package (v1.1.2) under a preference-space specification.
The mixing integral is approximated using 100 Sobol quasi-random draws
per respondent. To reduce sensitivity to starting values and mitigate
convergence to local optima, each model is initialized from 20
independent random starting points; the run achieving the globally
lowest log-likelihood is retained. Convergence is verified via the
relative and absolute function tolerance criterion (Exit Status 3) in
all six models.

Willingness to pay for attribute \(k\) is derived as the negative ratio
of the mean attribute coefficient to the mean price coefficient:

\[\widehat{\text{WTP}}_k = -\frac{\hat{\mu}_k}{\hat{\mu}_{\text{price}}}\]

This quantity represents the dollar amount the average respondent would
pay --- or require as a compensating discount --- for a one-unit change
in attribute \(k\), holding utility constant. Uncertainty bounds on WTP
are computed by drawing 10,000 samples from the estimated joint
distribution of \((\hat{\mu}_k,\, \hat{\mu}_{\text{price}})\) and
reporting the resulting 95\% simulation interval.

\begin{center}\rule{0.5\linewidth}{0.5pt}\end{center}

\subsection{4 Results}\label{results}

\subsubsection{4.1 WTP for used BEVs by adoption preference, vehicle
type, and driving
range}\label{wtp-for-used-bevs-by-adoption-preference-vehicle-type-and-driving-range}

Figure 1 presents WTP estimates at a fixed vehicle age of two years,
disaggregated by adoption preference (likely vs.~unlikely BEV adopter),
vehicle type (car vs.~SUV), and driving range (100, 200, and 300 miles).
The reference value of \$0 marks the ICE-parity threshold: WTP above
this line indicates a consumer-perceived premium for the used BEV
relative to an equivalent ICE vehicle; WTP below it indicates the
magnitude of discount required to reach purchase indifference.

\begin{figure}[H]

{\centering \includegraphics{vehicle_paper_2/images/dumbbell_plots/dumbbell_wtp_age_fixed_age_2_adopters.png}

}

\caption{Modelled WTP for used BEVs by adoption preference, vehicle
type, and range (vehicle age = 2 years). Points show point estimates;
horizontal lines show 95\% confidence intervals.}

\end{figure}%

Likely adopters consistently exhibit higher WTP than unlikely adopters
across all attribute combinations, and the gap is large. Among likely
car adopters at 300-mile range, WTP is the only estimate that crosses
into positive territory --- approximately +\$10,000 to +\$15,000 ---
indicating a willingness to pay a premium for a used BEV over an
equivalent ICE vehicle at this combination of attributes. All other
segments require some discount to induce purchase, with the required
discount increasing sharply as range declines and adoption propensity
decreases. Unlikely adopters at 100-mile range require discounts in the
range of \$35,000 to \$40,000 relative to an ICE alternative. Within
each adoption group, SUVs carry consistently lower WTP than cars at
equivalent range levels, a vehicle-type penalty that is stable across
the full range of estimates shown.

\subsubsection{4.2 WTP by adoption preference, vehicle type, and vehicle
age}\label{wtp-by-adoption-preference-vehicle-type-and-vehicle-age}

Figure 2 holds driving range constant at 300 miles and examines how WTP
varies across vehicle age cohorts (2, 5, and 8 years), disaggregated by
adoption preference and vehicle type.

\begin{figure}[H]

{\centering \includegraphics{vehicle_paper_2/images/dumbbell_plots/dumbbell_wtp_age_fixed_range_300_adopters.png}

}

\caption{Modelled WTP for used BEVs by adoption preference, vehicle
type, and vehicle age (range = 300 miles).}

\end{figure}%

Vehicle age exerts a large, negative, and monotonic effect on WTP. Each
successive age cohort shifts estimated WTP further below the parity
threshold, and the gradient is steeper than standard ICE depreciation
schedules would predict for the same age interval --- an excess discount
that most plausibly reflects consumer uncertainty about battery
state-of-health in older vehicles rather than mechanical wear per se.
The age penalty is most severe among unlikely adopters: 8-year-old SUVs
in this group require discounts approaching \$40,000. Likely adopters
are more resilient to vehicle aging --- their WTP for cars remains at or
near \$0 at five years of age --- but shifts meaningfully downward by
eight years. The joint variation in adoption preference and vehicle age
produces the widest spread of WTP estimates across the full set of model
outputs, establishing these two dimensions as the primary
value-determining axes in the used BEV market.

\subsubsection{4.3 WTP by budget tier, vehicle type, and vehicle
age}\label{wtp-by-budget-tier-vehicle-type-and-vehicle-age}

Figure 3 focuses on the high-budget consumer segment at 300-mile range
and presents WTP estimates across vehicle types and age cohorts.

\begin{figure}[H]

{\centering \includegraphics{vehicle_paper_2/images/dumbbell_plots/dumbbell_wtp_age_fixed_range_300_budgets.png}

}

\caption{Modelled WTP for used BEVs by vehicle type and age ---
high-budget consumers, range fixed at 300 miles.}

\end{figure}%

High-budget car buyers at two years of age represent the consumer
segment most proximate to price parity: the point estimate approaches
\$0 with a confidence interval that extends into positive territory,
indicating that near-parity pricing for used BEVs is a feasible market
outcome for this specific attribute combination. The age penalty
operates with similar force in this segment: the WTP differential
between 2- and 8-year-old vehicles is approximately \$15,000 to
\$20,000. Within the high-budget group, SUVs lag behind cars at every
age level, and confidence intervals are notably wider throughout this
segment relative to others, reflecting the greater estimated preference
heterogeneity among SUV buyers.

\subsubsection{4.4 Summary of attribute effects on
WTP}\label{summary-of-attribute-effects-on-wtp}

Table 1 summarizes the direction and relative magnitude of each
attribute's effect on WTP, which is consistent in sign across all model
specifications.

\begin{longtable}[]{@{}lll@{}}
\toprule\noalign{}
Attribute & Direction of effect on WTP & Relative magnitude \\
\midrule\noalign{}
\endhead
\bottomrule\noalign{}
\endlastfoot
Driving range (100 → 300 miles) & Positive & Large \\
Vehicle age (2 → 8 years) & Negative & Large \\
Adoption preference (unlikely → likely) & Positive & Large \\
Budget tier (low → high) & Positive & Moderate \\
Vehicle type (SUV vs.~car) & Negative for SUV & Moderate \\
\end{longtable}

\textbf{Table 1.} Direction and relative magnitude of attribute effects
on WTP for used BEVs across all model specifications.

Driving range, vehicle age, and adoption preference are the dominant
value determinants, each producing effects that substantially exceed
those of budget tier and vehicle type. The only segment for which WTP
reaches or exceeds the ICE-parity threshold is likely BEV adopters
selecting 300-mile-range cars at two years of age. For all other
combinations, used BEVs must be offered at a discount relative to
equivalent ICE vehicles to attract buyers, with the required discount
scaling predictably with increasing vehicle age, declining range, and
greater adoption reluctance.

\subsubsection{4.5 Mixed Logit Estimates: Preference Heterogeneity by
Adoption
Group}\label{mixed-logit-estimates-preference-heterogeneity-by-adoption-group}

The six MXL models establish that the aggregate WTP patterns in Figures
1--3 reflect qualitatively distinct preference structures across
adoption groups, and that substantial within-group heterogeneity
characterizes both segments. Among likely BEV adopters, mean BEV
powertrain coefficients are negative across all three vehicle-segment
specifications (full sample: β = −0.49; car: β = −1.01; SUV: β = −1.88),
yet are in each case accompanied by large and statistically significant
standard deviations (\textbar σ\textbar{} = 2.2--3.2, p \textless{}
0.001 throughout). The apparent contradiction between negative mean BEV
utility and self-reported adoption likelihood is resolved by the mixing
distribution: simulation-based draws from the estimated parameter
distributions confirm that the interquartile range for BEV utility
straddles zero within each subgroup, indicating genuine within-group
polarization in BEV attitudes rather than a uniformly receptive
population. HEV powertrains are, by contrast, unambiguously and
consistently preferred over the ICE baseline among likely adopters (β =
0.89--0.95, p \textless{} 0.001), suggesting that hybrid electrification
represents the more broadly accepted form of powertrain transition
within this group. Driving range is the dominant positive attribute
among likely car adopters specifically (β = 1.061, p \textless{} 0.001)
--- approximately 2.5 times its magnitude in the full likely-adopter
sample --- identifying range as the primary attribute lever for this
subgroup. McFadden R² values across the three likely-adopter models
range from 0.164 to 0.214, consistent with good fit for discrete choice
models; the SUV segment yields the lowest fit (R² = 0.164), reflecting
greater idiosyncratic preference variation among SUV buyers who are
otherwise positively oriented toward BEVs.

Among unlikely BEV adopters, the preference structure is categorically
different in both direction and magnitude. Mean BEV utility coefficients
are large and negative across all three models (full sample: β = −7.78;
car: β = −9.27; SUV: β = −8.35; all p \textless{} 0.001), with
simulation-based draws confirming near-universal BEV aversion in the car
segment: at least 75\% of the simulated population holds strongly
negative BEV utility, with effectively no probability mass above zero.
HEV powertrains are also negatively valued in this group (β = −0.62 to
−0.76, p \textless{} 0.001), establishing that the resistance is to
electrification broadly rather than to BEVs specifically. Price
sensitivity among unlikely adopters substantially exceeds that of likely
adopters: the price coefficient for car-segment unlikely adopters (β =
−2.912) is 63\% larger in magnitude than for the corresponding
likely-adopter group (β = −1.792), and the full-sample differential is
approximately 29\%. Driving range is a statistically significant
predictor of vehicle choice in five of the six models; the single
exception is the car-segment unlikely-adopter model (β = 0.420, p =
0.272), which is the only model in which range fails to reach
significance and indicates that range improvements alone would not
induce BEV consideration among this segment's car buyers. Despite the
pervasive aversion, large standard deviations on both BEV and HEV
coefficients (\textbar σ\_BEV\textbar{} = 5.2--6.5;
\textbar σ\_HEV\textbar{} = 2.7--3.4, p \textless{} 0.01 in all cases)
confirm within-group heterogeneity sufficient to indicate that a
non-trivial minority of unlikely adopters may be reachable through
targeted interventions --- particularly in the SUV segment, where mean
HEV aversion is weakest (β = −0.617).

\begin{center}\rule{0.5\linewidth}{0.5pt}\end{center}

\subsection{5 Discussion and Policy
Implications}\label{discussion-and-policy-implications}

\subsubsection{5.1 Pricing and market
structure}\label{pricing-and-market-structure}

The results collectively establish that the used BEV market has not
reached price parity with conventional vehicles for the majority of
prospective buyers. Only the intersection of likely adoption propensity,
300-mile driving range, car vehicle type, and low vehicle age produces
WTP estimates at or near the ICE-equivalent price; every other segment
combination requires a discount of varying magnitude. This finding has
direct implications for dealer pricing and remarketing strategy: used
BEVs aged beyond approximately five years, or those with rated ranges
below 200 miles, face a valuation gap that cannot be closed by marginal
price concessions alone.

The age effect on WTP warrants particular attention. The estimated WTP
differential between 2-year-old and 8-year-old vehicles --- on the order
of \$15,000 to \$20,000 within a given segment --- considerably exceeds
what standard ICE depreciation schedules predict for the same age
interval, and most plausibly reflects consumer uncertainty about battery
state-of-health rather than mechanical depreciation per se.
Interventions that credibly reduce SOH uncertainty --- standardized
battery health disclosure requirements, third-party certification
programs, or manufacturer-extended battery warranties in the secondary
market --- could recover a portion of this excess age discount without
requiring changes in the underlying vehicle. In the absence of such
mechanisms, downward pressure on used BEV residual values will intensify
as early-adopter vehicles age into the 5--8 year range in coming years.

\subsubsection{5.2 Segment-differentiated intervention
strategies}\label{segment-differentiated-intervention-strategies}

The breadth of the WTP distribution across adoption groups indicates
that no single policy or market instrument is uniformly effective. Among
likely BEV adopters --- the population already proximate to the purchase
threshold --- the binding constraints are driving range and total cost
of ownership rather than powertrain aversion. The MXL results identify
range as the largest positive attribute in the car-segment model and
operating cost as the dominant financial disutility factor in that
group, with the operating cost coefficient exceeding the price
coefficient. These estimates point to consumer communication strategies
emphasizing per-mile cost advantages and range sufficiency for typical
daily use as the most cost-efficient mechanism for converting this
population. The strong opt-out aversion observed among likely adopters
further confirms high baseline purchase engagement, reducing the
intervention threshold for this segment.

Among unlikely BEV adopters, the constraints are more fundamental. BEV
aversion in this population is not contingent on range adequacy --- the
non-significant range coefficient in the car-segment model directly
excludes range as a primary driver --- and price sensitivity is 29--63\%
higher than among likely adopters, implying that financial incentives
would need to be substantially larger than current programs provide to
shift revealed preferences at scale. The most tractable near-term lever
identified by the data is the significant within-group heterogeneity on
HEV coefficients: the large standard deviations and the relatively mild
mean HEV aversion among SUV-segment unlikely adopters (β = −0.617)
identify hybrid SUVs as a more accessible intermediate adoption pathway
than direct BEV promotion for this population.

\subsubsection{5.3 Implications for incentive
design}\label{implications-for-incentive-design}

The federal used EV tax credit, currently capped at \$4,000 for
qualifying used vehicle transactions, is substantially smaller than the
price concessions required to bring most unlikely-adopter segments to
purchase indifference. Required discounts for unlikely adopters at
realistic age and range combinations range from approximately \$15,000
to over \$35,000, indicating that the current federal instrument is
insufficient to materially expand used BEV uptake beyond buyers who
would likely have purchased regardless. Achieving broader market
penetration would require either substantially larger per-vehicle
transfers or complementary non-price interventions --- charging
infrastructure deployment, battery health certification, consumer
education --- that reduce the risk premium currently embedded in used
BEV valuations for the majority of prospective buyers.

State-level programs designed with segment awareness may achieve better
marginal returns by concentrating resources on consumers closest to the
purchase threshold. Targeting high-budget car buyers at low vehicle age
and long range --- where WTP already approaches parity --- minimizes the
transfer required per incremental transaction. Uniform subsidies applied
without regard to vehicle age, range, or buyer propensity risk
distributing a substantial share of resources to inframarginal
transactions and generating limited additionality.

\subsubsection{5.4 Limitations and future
directions}\label{limitations-and-future-directions}

Several limitations bound the scope and external validity of these
findings. First, the analysis rests on stated preferences, which may
diverge from revealed purchase behavior; hypothetical bias is a
well-documented concern in DCE research, particularly for product
categories with which respondents lack direct experience, and the
magnitude and direction of this bias in the used BEV context is
uncertain. Second, the analytic sample was recruited through online
research panels, which may overrepresent digitally engaged and
information-seeking consumers relative to the broader used-vehicle
buying population. Third, the MXL specification models preference
heterogeneity as continuous normal distributions; this may inadequately
represent the possibility of discrete, internally coherent preference
clusters within adoption groups, and a latent class analysis would test
that possibility directly. Fourth, battery state-of-health --- arguably
the most consequential unobserved attribute in the used BEV transaction
--- was not varied in the vehicle powertrain DCE analyzed here; its
effect on WTP is addressed in a companion battery DCE component of the
broader study. Integrating findings from both DCE components and
cross-validating WTP estimates against revealed transaction prices from
dealer records would substantially strengthen the empirical basis for
the inferences drawn above.

\begin{center}\rule{0.5\linewidth}{0.5pt}\end{center}

\subsection{References}\label{references}

{[}1{]} Bureau of Transportation Statistics (2023). \emph{New and Used
Passenger Car and Light Truck Sales and Leases}. www.bts.gov.

{[}2{]} Cox Automotive Inc.~(2025). \emph{February 2025 Used EV Sales
Data}.

{[}3{]} Hagman, J., Ritzén, S., Stier, J. J., \& Susilo, Y. (2016).
Total cost of ownership and its potential implications for battery
electric vehicle diffusion. \emph{Research in Transportation Business \&
Management}, 18, 11--17.

{[}4{]} Fernandes, M. (2023). \emph{Consumer survey on used EV
consideration}.

{[}5{]} Prieto, M., Caemmerer, B., \& Baltas, G. (2015). Fuel type
choice in the used car market. \emph{Transportation Research Part A},
80, 1--15.

{[}6{]} Consumer Reports; U.S. EPA (various years). \emph{Used Car
Buying Guides}.

{[}7{]} Kelley Blue Book (2025, April). \emph{Used Car Market Report}.

{[}8{]} Tandf (2025). \emph{Consumer segments in the used EV market}.

{[}9{]} Sd (2024). \emph{Demographics of used EV owners}. Studies 1 \&
2.

{[}10{]} Trd (2024). \emph{Pre-purchase concerns and income sensitivity
in used EV adoption}.

{[}11{]} Dnistran (2024). \emph{Battery degradation analysis across
\textasciitilde5,000 EVs}.

{[}12{]} Sheykhfard, A. et al.~(2025). \emph{Used EV ownership
experience and battery concerns}.




\end{document}
